\documentclass[../principal.tex]{subfiles}
\graphicspath{{\subfix{../imagenes/}}}

\begin{document}

\chapter{Funciones}

Las funciones son bloques de código que realizan una tarea. Las funciones tienen un comportamiento plástico y flexible gracias al uso de parámetros.

La sintaxis de las funciones en C++ es la siguiente:

\begin{lstlisting}[language=C++]
tipo nombre(tipo unParam, tipo otroParam) {
}
\end{lstlisting}

Si has programado un microcontrolador Arduino como la placa Uno R4 Minima, el código de inicio muestra la siguiente configuración

\begin{lstlisting}[language=C++]
void setup() {

}

void loop() {

}
\end{lstlisting}


En una placa del tipo Raspberry Pi Pico, por ejemplo la 2W, encontrarás

\begin{lstlisting}[language=C++]
    
int main() {

}
\end{lstlisting}


\newpage

\end{document}

